% Copyright (c) 2026 by C. Klukas.
% Licensed under the MIT License. See LICENSE for details.

% This appendix is dense with monospaced option and script names that
% cannot be hyphenated; the emergency stretch keeps such paragraphs inside
% the text width (as in Appendix A).
\begingroup\setlength{\emergencystretch}{3em}

Every number in Section~\ref{sec:evaluation}, every table of recorded
results and every code excerpt whose caption names a source path is
produced by a script or a tool in the repository \cite{pagletsCpp}. This
appendix names them, in the order in which one would run them: build the
edition, run the test suite and record its results, run the live
measurements on real machines, run the protocol checker, and regenerate the
documentation and this paper from the recorded results.
Figure~\ref{fig:reproduce-pipeline} shows how the pieces connect. The
commands below run from the repository root of the branch \code{cpp}; those
that concern this paper (Section~\ref{sec:reproduce-generated}) need the
branch \code{cpp-paper}, which carries \code{papers/} and takes \code{cpp}
in by merge. The page \code{building.md} of the user documentation
\cite{pagletsCppDocs} describes the build and the test suite in more
detail; \code{examples/index.md} (``Test results'') says how the examples'
results are recorded.

\begin{figure}[htbp]
\centering
\begin{tikzpicture}[node distance=5mm and 4mm]
  \node[box] (tests) {\code{paglets\_tests}\\\footnotesize\code{PAGLETS\_RESULTS\_DIR}};
  \node[box, right=of tests] (live) {\code{mesh\_bench.sh}\\\code{chronicle\_live.sh}};
  \node[box, right=of live] (check) {\code{paglets-protocol-check}\\\footnotesize\code{-{}-results}};
  \node[plain, below=of live] (json) {\code{docs/results/*.json}};
  \node[box, left=of json, xshift=-6mm] (src) {sources\\\footnotesize\code{// -{}-8<-{}- [start:NAME]}};
  \node[state, below=of json] (sync) {\code{paglets-docs-sync}};
  \node[plain, below left=of sync, xshift=-2mm] (md) {\code{docs/*.md}};
  \node[plain, below right=of sync, xshift=2mm] (tex) {\code{papers/.../chapters/*.tex}};
  \node[box, below=of md] (site) {\code{ckdocs}: the documentation site};
  \node[box, below=of tex] (pdf) {\code{latexmk}, \code{gnuplot}: this paper};
  \draw[arrow] (tests) -- (json);
  \draw[arrow] (live) -- (json);
  \draw[arrow] (check) -- (json);
  \draw[arrow] (json) -- (sync);
  \draw[arrow] (src) |- (sync);
  \draw[arrow] (sync) -- (md);
  \draw[arrow] (sync) -- (tex);
  \draw[arrow] (md) -- (site);
  \draw[arrow] (tex) -- (pdf);
\end{tikzpicture}
\caption{Where the numbers and the excerpts come from. Tests, live scripts
and the protocol checker record results as JSON; \code{paglets-docs-sync}
writes tables from them, and code excerpts from marked regions of the
sources, into the documentation pages and into the chapters of this paper.}
\label{fig:reproduce-pipeline}
\end{figure}

\subsection{Building}
\label{sec:reproduce-building}

The host is C++26 code built with GCC~16 (\code{-std=c++26 -freflection});
reflection generates the wire codecs and the guest schema code. Guest
modules are compiled by clang for \code{wasm32-wasip1} with a WASI sysroot.
The build fetches WAMR and cpp-httplib at pinned releases. On macOS the
toolchain comes from Homebrew, on Linux GCC~16 is combined with the
wasi-sdk in \code{/opt/wasi-sdk} (or \code{WASI\_SDK\_PATH}):

\begin{lstlisting}[style=paglets-bash]
# macOS
brew install gcc@16 llvm lld wasi-libc wasi-runtimes cmake ninja \
    libsodium openssl@3 zstd pkgconf
# Linux
apt install libsodium-dev libssl-dev libzstd-dev pkg-config
\end{lstlisting}

The build and the tests use CMake presets (Table~\ref{tab:reproduce-presets}).
The presets configure \code{RelWithDebInfo} builds; the results recorded in
\code{docs/results/} come from Release builds (configured with
\code{CMAKE\_BUILD\_TYPE} set to \code{Release}), as their \code{build}
field says:

\begin{lstlisting}[style=paglets-bash]
cd cpp
cmake --preset macos-arm64
cmake --build --preset macos-arm64
ctest --preset macos-arm64
build/macos-arm64/host/paglets-host --info
\end{lstlisting}

\code{paglets-host -{}-info} prints what a build contains: version,
platform, compiler, whether reflection is on, the WAMR release with its
interpreter and bound-check mode (for example \code{fast-interp, hardware
bound checks}), the paglet ABI and the mesh protocol. The build directory
is \code{cpp/build/<preset>}; the guest modules end up in its
\code{guests/} directory. Two configure options matter for the
measurements. \code{PAGLETS\_WAMR\_FAST\_INTERP}, on by default, selects
WAMR's fast interpreter. \code{PAGLETS\_WAMR\_HW\_BOUND\_CHECK} must be
off on kernels with a 39-bit address space such as Raspberry Pi OS, where
guard-page bound checks limit a process to about 60 instances. A Linux
host of this kind (a GCC~16 installed next to an older system compiler, as
on Raspberry Pi OS) is built with the recipe in \code{docs/building.md}
(``Build and test''): a Release build of \code{paglets-host} and
\code{paglets-worker} only, guests and tests off, hardware bound checks
off, the C++ library linked statically. Such a host needs no guest modules
of its own; the modules travel from the machine that launches the paglets.
The recorded runs describe the hosts only as
Table~\ref{tab:eval-hosts} shows them; how each host was built is
not recorded.

\begin{table}[htbp]
\centering\small
\caption{The CMake presets (\code{cpp/CMakePresets.json}).}
\label{tab:reproduce-presets}
\begin{tabular}{ll}
\toprule
Preset & Use \\
\midrule
\code{macos-arm64}, \code{linux-gcc16} & the regular builds \\
\code{windows-mingw-gcc16} & the regular build with MinGW-w64 (MSYS2 UCRT64) \\
\code{macos-arm64-asan}, \code{linux-gcc16-asan} & AddressSanitizer, UndefinedBehaviorSanitizer \\
\code{linux-gcc16-tsan} & ThreadSanitizer \\
\code{linux-host-only} & a C++23 compiler without reflection; prebuilt guests \\
\bottomrule
\end{tabular}
\end{table}

\subsection{The test suite and the recorded results}
\label{sec:reproduce-tests}

\code{ctest} runs \code{paglets\_tests} twice, with \paglet{}s inside the
test process (\code{unit}) and in sandboxed worker processes
(\code{unit\_workers}), and a dozen tests of the programs: the sandbox
check, memory images saved and resumed across processes, hosts that serve,
move and find each other, the CLI ledger, and the installed package
(\code{docs/building.md}, ``The test suite''). The \code{unit} test also
explores six small worlds of the zombie protocol exhaustively and runs 200
timed random runs (\code{cpp/tests/test\_protocol\_check.cpp}). A further
test, \code{docs\_sync}, fails while a documentation page or a chapter of
this paper is not current (Section~\ref{sec:reproduce-generated}).

The tests of the example paglets record what they measure. With
\code{PAGLETS\_RESULTS\_DIR} set, a test that reports to a topic is written to
\code{<dir>/<topic>.json} after it ran: its name, mode, outcome, duration,
platform, build type, date and the values it reported
(\code{cpp/tests/test.hpp}). One script runs every recording test, in
process and in worker processes, then the full protocol check, and then
the generator:

\begin{lstlisting}[style=paglets-bash]
cpp/tools/update_results.sh [cpp/build/macos-arm64]   # the build directory (default)
\end{lstlisting}

It runs \code{paglets\_tests} with the filters \code{example:},
\code{gateway:}, \code{exchange:} and \code{pi:}, the hello and ping-pong
samples and the counter snapshot test, each once in the test process and
once in worker processes (\code{PAGLETS\_TEST\_WORKER}). Then it runs the
protocol checker with the full suite and 10\,000 timed runs
(\code{-{}-suite full}, \code{-{}-simulate}, \code{-{}-seeds 10000}) and
records the outcome in \code{zombie\_protocol.json}; the recorded run
(\code{docs/results/zombie\_protocol.json}: Darwin arm64, Release build,
2026-10-13) took 2393\,s in all, about 40 minutes, nearly all of it in the
exhaustive worlds; the 10\,000 timed runs took 28\,s. Finally it runs
\code{paglets-docs-sync}. The results are committed under
\code{docs/results/}, one file per topic.

\subsection{The live measurements}
\label{sec:reproduce-live}

Three scripts run the examples on real machines: the machine that runs the
script (\emph{this machine} below; the Mac in the recorded runs) and remote
Linux hosts reached over SSH. Each builds a fresh mesh for the run (keys,
a ledger with one admin and one owner, every host enrolled), starts
\code{paglets-host serve} on every host with explicit \code{-{}-peer}
addresses and beacons off, waits until \code{remote landscape} shows every
host, and tears everything down at the end (\code{chronicle\_live.sh
-{}-keep} leaves the mesh running for a look). Every host needs
\code{paglets-host} and \code{paglets-worker} built from the same sources;
the modules travel from this machine. Table~\ref{tab:reproduce-scripts}
lists the scripts, their options are in their headers, and
Table~\ref{tab:eval-hosts} shows the machines of the recorded runs as
\code{mesh\_bench.sh} describes them (the Mac launches the paglets; the
Linux hosts run only the host and the worker).

\begin{lstlisting}[style=paglets-bash]
cpp/tools/mesh_bench.sh \
    --remote pi=klukas@rpi4:/home/klukas/paglets-live/cpp/build/linux-gcc16/host \
    --remote uconsole=uconsole:/home/klukas/paglets-live/cpp/build/linux-gcc16/host \
    [--local-name mac] [--build cpp/build/macos-arm64] [--port 7443] \
    [--sizes "16 64 256"] [--pi-digits 100000] \
    [--steps "network transfer latency benchmark pi"]
cpp/tools/chronicle_live.sh --remote pi=... --remote uconsole=... \
    [--lifetime-min 10] [--no-failures] [--keep] [--late NAME]
cpp/tools/web_live.sh --remote pi=... [--check-only]
\end{lstlisting}

\code{mesh\_bench.sh} runs five steps. Four write a results file of their
own (\code{live\_transfer}, \code{live\_latency}, \code{live\_benchmark},
\code{live\_pi}); the network step's numbers go into every transfer entry
as its TCP baseline, and \code{live\_hosts} is written on every run.
\code{-{}-steps} selects steps; the files of the others stay as they are.
\emph{network}: iperf3 over TCP for 5 seconds between every ordered pair
of hosts. \emph{transfer}: a File Courier carries a file of random bytes of
every size in \code{-{}-sizes} (MB) for every ordered pair; the rate is the
size divided by the time the destination took to save it (fetching every
part, checking the SHA-256, writing). \emph{latency}: the Latency Map, 20
round trips between paglets for every ordered pair of hosts.
\emph{benchmark}: the Mesh Benchmark in a compute slot on every host, with a
budget of one second, 4\,MB of memory and 8\,MB of storage. \emph{pi}:
\code{-{}-pi-digits} hex digits of pi in chunks of 500, up to 32 in flight,
streamed in order to a file on this machine and converted to decimal; the
step passes when the computation ends in the state \code{done} and the
first 100 decimal digits are right. The hosts' logs are kept in
\code{\$TMPDIR/paglets-bench-logs} for a look afterwards.

\code{chronicle\_live.sh} needs at least two remote hosts. It starts one
Mesh Chronicle with a lifetime of ten minutes (\code{-{}-lifetime-min})
whose life log goes to this machine, and injects failures at fixed
fractions of the lifetime: the host of the current main is frozen with
\code{SIGSTOP} at 15\,\% and thawed at 30\,\%, so that two mains meet; the
host of the main is killed at 45\,\% and restarted with its state at
60\,\%; another host is killed at 70\,\% and restarted at 85\,\%; the
sink's host shuts down normally at 87.5\,\% and is back at 92.5\,\%. With
\code{-{}-late NAME} that host joins the mesh only at 10\,\%. Afterwards
the script fetches the life log with \code{remote zombies result}, runs
\code{remote zombies check}, which checks the protocol steps every live
host recorded against the specification, and counts the replicas left. It
records the outcome in \code{docs/results/live\_chronicle.json} and the
life log, as delivered, in \code{docs/examples/mesh-chronicle-log.md}.

\code{web\_live.sh} records nothing: it puts a counter on every host, a
\code{hello} and a Mesh Chronicle on the same kind of mesh, checks the API
of \code{paglets-web} against it (every host reports its status; a call
through the UI reaches a \paglet{} on a remote host), and leaves the mesh
running for a look at the page (\code{-{}-check-only} stops after the
checks).

\begin{table}[htbp]
\centering\small
\caption{The live scripts in \code{cpp/tools/}, what they need on this machine beyond the build, and what they record.}
\label{tab:reproduce-scripts}
\begin{tabular}{lll}
\toprule
Script & Needs & Records in \code{docs/results/} \\
\midrule
\code{mesh\_bench.sh} & jq, iperf3 (on every host) & \code{live\_hosts}, \code{live\_transfer}, \code{live\_latency}, \\
 & & \code{live\_benchmark}, \code{live\_pi} \\
\code{chronicle\_live.sh} & jq & \code{live\_chronicle}; the life log in \code{docs/examples/} \\
\code{web\_live.sh} & jq, curl & nothing: a check of the web UI on a real mesh \\
\code{update\_results.sh} & a build with tests & one topic per example, and \code{zombie\_protocol} \\
\bottomrule
\end{tabular}
\end{table}

\subsection{The protocol checker}
\label{sec:reproduce-checker}

\code{paglets-protocol-check} (\code{cpp/tools/protocol\_check}) checks the
zombie protocol of Section~\ref{sec:zombies} in the two ways described
there: every reachable state of a small world, and timed random runs with
virtual clocks. Both run the same transition function that the guest SDK
runs.

\begin{lstlisting}[style=paglets-bash]
C=cpp/build/macos-arm64/tools/protocol_check/paglets-protocol-check
$C --suite full                       # the twelve worlds of the documentation
$C --simulate --seeds 10000           # timed random runs (default: 1000)
$C --trace 42                         # one simulated run: role changes and faults
$C --replicas 2 --crashes 1 --cuts 1  # one world, given by its options;
                                      # options not given keep the model's defaults
\end{lstlisting}

The suite \code{quick} is the set of worlds the \code{unit} test explores:
two replicas with one kind of fault each, and two replicas with a host
joining later and a crash. The suite \code{full} is the set recorded in
the documentation and in Section~\ref{sec:evaluation}: ten two-replica
worlds with two or three faults or owner actions (a cancel, a redirect)
each, explored completely, and two three-replica worlds with at most two
messages in flight per channel, explored to a depth of eleven steps, for
safety only. A world can also be given by its options (the usage line of
the tool lists them: replicas, work items, observations, each kind of
fault, messages in flight per channel, a state limit). Options that are not
given keep the defaults of \code{ModelConfig}
(\code{cpp/common/include/paglets/protocols/zombie/model.hpp}): three
replicas, one crash, cut, loss, duplicate and early timeout each, three
messages in flight; the suites start from no faults instead. A violated
property is reported with the shortest sequence of steps that leads to it.
The exit status is 0 when every property holds, 1 when one is violated
(or, for the full suite, when a role change of the table was reached by no
world), and 2 for a usage error. \code{-{}-results FILE} writes the outcome
in the format of \code{docs/results/}, which is how
\code{update\_results.sh} records it; \code{-{}-quiet} reports a violated
property without printing its steps.

\subsection{Generated documentation and this paper}
\label{sec:reproduce-generated}

The markers that \code{paglets-docs-sync} expands (Section~\ref{sec:tools})
have a LaTeX form in the chapters of this paper: they are LaTeX comments,
\code{\% snippet: PATH\#SECTION} or \code{\% results: TOPIC} followed by
listings options, and the generator writes a \code{lstlisting} (with the
source path appended to the caption, when the options give one), or a
booktabs table. Every listing in this paper whose caption ends in a source
path, and every table of recorded results
(Table~\ref{tab:eval-hosts} is one), was generated this way; the
\code{docs\_sync} test fails while any of them is out of date. The
generator needs only a C++23 compiler; the paper's workflow compiles it
with one \code{g++} command. The documentation site is built with ckdocs
from \code{docs/} and checked with \code{ckdocs check}.

\begin{lstlisting}[style=paglets-bash]
S=cpp/build/macos-arm64/tools/docs_sync/paglets-docs-sync
$S --root .            # rewrite docs/ (and papers/, on cpp-paper)
$S --root . --check    # change nothing; exit 1 if a page or a chapter is stale
cd papers/paglets-cpp  # on the branch cpp-paper
make sync              # the same, with the build in PAGLETS_BUILD
make check
make                   # the PDF: figures/data.sh (jq), gnuplot, latexmk
\end{lstlisting}

gnuplot draws the figures of Section~\ref{sec:evaluation} from TSV files
in \code{build/data/}, which \code{figures/data.sh} extracts from the same
JSON files with jq. The paper is built by a GitHub workflow on every push
to the branch \code{cpp-paper} that changes the paper, the recorded results
or \code{cpp/}: it checks that the excerpts and tables are current, builds
the PDF, publishes it to a release tagged \code{cpp-paper-YYYY-MM-DD}, and
starts the documentation workflow, which places the newest PDF on the
documentation site.

\subsection{Installing with Homebrew}
\label{sec:reproduce-homebrew}

On macOS the tap \code{cklukas/homebrew-paglets} builds the latest release
from source. The formula depends on Homebrew's GCC, because its build turns
reflection on and clang has none (\code{fails\_with :clang}), and on
libsodium, OpenSSL~3 and zstd; it installs \code{paglets-host},
\code{paglets-worker}, \code{paglets-web}, the guest SDK and the CMake
package for \code{find\_package(paglets)}. Building paglets against it
needs the guest toolchain as well.

\begin{lstlisting}[style=paglets-bash]
brew install cklukas/paglets/paglets           # the latest release (tag cpp-vX.Y.Z)
brew install --HEAD cklukas/paglets/paglets    # the current cpp branch
brew install llvm lld wasi-libc wasi-runtimes  # the guest toolchain
\end{lstlisting}

The installed package builds guest modules and runs them with
\code{paglets-host}; it does not export the host runtime as a library.
\par\endgroup
